Test-time training lets a language model improve on a single hard
problem at inference time, but on-policy self-distillation stalls when
the model rarely produces a correct attempt to learn from. This thesis
studies how to accelerate test-time discovery in complex code generation
via \emph{execution-guided self-distillation}, where execution feedback
(failing tests, runtime errors) is the privileged context that drives
the update. We propose a \emph{teacher-first} organization: rather than
distilling the student's own failed rollout, the self-teacher generates
trajectories under feedback, a verifier and a judge filter them for
correctness and independence, and the student is distilled toward those
filtered trajectories, placing the method between on-policy SDPO and
off-policy SFT. Across a medium-sized matched evaluation on hard
LiveCodeBench problems with a 4B model, teacher-first significantly
outperforms the student-first baseline (Wilcoxon signed-rank), and on
the hardest problems it escapes the flat-reward trap that keeps
student-first stuck at zero. A compute-matched comparison shows the
advantage is not a sampling-volume artifact: at equal generation budget
teacher-first still wins, and reaches a given discovery probability with
fewer total generations. With thinking enabled, we measure epistemic
verbalization at test time and find that distillation from
answer-bearing context suppresses uncertainty markers, connecting the
method to the degradation Kim et al.~describe. Finally, supplying a
method-bearing reference (worked solutions from MATH-500) lets the model
escape on math problems where an answer-only reference does not,
validating the central \emph{value-versus-procedure} boundary: escape
requires the reference to encode an in-reach method, not merely a value.

\vspace{0.5cm}
\textbf{Keywords:} test-time training, self-distillation, code
generation, execution feedback, discovery@k, epistemic suppression,
compute-to-correct.